\section{Introduction}\label{sec:intro}

The employment effects of the minimum wage remain among the most debated questions
in labor economics, and the debate is sharper still in developing countries, where a
large share of the workforce operates outside the reach of labor regulation. In such
settings a statutory wage floor can raise pay for covered workers, push some of them
into the uncovered informal sector, or serve as a coordinating reference for wages
that formal firms would have set in any case. Which of these forces dominates is an
empirical question whose answer depends on how binding the floor is, how strictly it
is enforced, and how easily workers and firms move between the formal and informal
margins. This paper studies these questions for Peru, a country that combines a high
statutory minimum wage relative to median pay with one of the highest rates of labor
informality in Latin America.

We exploit the increase in Peru's national minimum wage (the \emph{Remuneraci\'on
M\'inima Vital}, or RMV) from 930 to 1,025 soles, enacted by Supreme Decree
003-2022-TR and effective on 1 May 2022. This 10.2 percent nominal increase was the
first change to the RMV in four years, and it occurred as the labor market was
emerging from the pandemic. A central institutional fact shapes the analysis: Peru
sets a single national minimum wage, with no regional or sectoral variation in the
general regime. There is therefore no cross-sectional variation in the level of the
policy and no staggered adoption to exploit. As \citet{jimenez_2023} notes for the
2016 reform, credible identification must instead come from differential
\emph{exposure} to a policy that changes for everyone at the same time.

Our source of differential exposure is education. In the pre-reform quarters the
median monthly wage of a dependent employee in Peru was almost exactly 930 soles,
the old minimum, so the typical low-education worker sat right at the wage floor. A
10 percent increase in that floor is highly binding for such workers. High-education
workers, by contrast, have a median wage well above the minimum and are far less likely to
be bound by its level, although some remain exposed through informal or own-account work.
This gap in exposure, rather than a claim of zero exposure for the comparison group,
motivates a difference-in-differences design that
contrasts low-education workers (those with at most complete secondary schooling,
who make up roughly two thirds of the working-age population) with high-education
workers (those with any tertiary schooling). The design follows a long tradition of
measuring minimum wage effects from variation in the ``bite'' of the policy across
groups or regions \citep{card_1992, lee_1999, machin_manning_rahman_2003}, adapted
here to a setting with a single national reform and rich, quarterly household data.

We use the public-use quarterly employment modules of the Encuesta Nacional de
Hogares (ENAHO) from the first quarter of 2021 through the fourth quarter of 2024,
sixteen quarters that bracket the reform with five pre-reform and eleven post-reform
periods. The panel structure of the data at the quarterly frequency lets us estimate
a dynamic event study, test for differential pre-trends by skill, and trace the path
of effects for more than two years after the reform. We examine a broad set of
outcomes: real hourly wages and monthly earnings, the probability of employment and
of labor force participation, the incidence of formal versus informal employment,
hours of work, and the composition of employment between wage work and
self-employment. Throughout we use survey weights and cluster inference by
department, and we complement conventional cluster-robust standard errors with the
small-sample correction of \citet{pustejovsky_tipton_2018} and with randomization
inference, because the treatment is assigned at a coarse level.

Three findings emerge. First, the reform did not raise the real wages of
low-education workers relative to high-education workers. The estimated effect on the
log real hourly wage is a precisely estimated null, and the event study shows no
break at the reform and flat pre-trends. This is our most robust result. Second, the
reform is associated with a decline in the probability that a low-education worker
holds formal salaried employment and with an offsetting rise in self-employment, a
pattern consistent with reallocation from the formal to the informal margin rather
than with the compression of the formal wage distribution documented for earlier
Peruvian reforms. Third, the employment rate of low-education workers falls modestly
relative to that of high-education workers, though this margin is the least robust
and is sensitive to the treatment of the pandemic recovery quarter of early 2021.
For the wage and formality results, pre-trends are statistically flat; for
employment and self-employment they are not, and we address this directly using the
sensitivity analysis of \citet{rambachan_roth_2023} and a battery of specification
and placebo checks. We further replicate the entire design on the subsequent reform
of January 2025, which raised the RMV to 1,130 soles in a cleaner, post-pandemic
window, and obtain a consistent picture.

Our contribution is threefold, and we are careful not to overstate it. First, we
provide the first evaluation of the 2022 Peruvian minimum wage reform and, to our
knowledge, the first study of any Peruvian reform using the post-pandemic quarterly
ENAHO. The existing micro evidence for Peru stops at the 2016 reform
\citep{jimenez_2023} or the reforms of the early 2000s \citep{jaramillo_lopez_2006,
cespedes_2005}. Second, we make heterogeneity by education the central object of
analysis at the worker level. Prior work cuts heterogeneity by income range or age,
or constructs a department-level ``dose'' of exposure; none contrasts low- and
high-education workers directly as treatment and comparison groups in an event study
of the Peruvian minimum wage. Third, we place the informal margin at the center of
the analysis and show that the adjustment to the reform runs through the composition
of employment rather than through pay. We do not claim novelty for the underlying
identification idea, which is the standard exposure or bite design, nor for the
econometric toolkit, which is by now well established. The value added is the
combination of a new reform, a new and policy-relevant period, an education-based
worker-level design, and a careful treatment of the informal margin and of the
threats to identification that such a design entails.

The remainder of the paper is organized as follows. Section~\ref{sec:institutional}
describes the Peruvian minimum wage and the reform. Section~\ref{sec:literature}
reviews the related literature. Section~\ref{sec:data} presents the data and the
construction of the key variables. Section~\ref{sec:strategy} sets out the empirical
strategy and the identifying assumptions. Section~\ref{sec:results} reports the main
results, Section~\ref{sec:robustness} the robustness and inference checks, and
Section~\ref{sec:heterogeneity} the heterogeneity and distributional analysis.
Section~\ref{sec:discussion} discusses mechanisms, external validity, and
limitations, and Section~\ref{sec:conclusion} concludes.
